\chapter{Literature Review}
\label{ch:litreview}

\section{Critical Assessment of Prior Works}
We synthesise and critically analyse the current state of student simulation and Socratic tutoring, highlighting achievements, trajectories, omissions, and how our work addresses these gaps.

\subsection{Epistemic Fidelity vs. Surface Realism}
Yuan et al. \cite{yuan2026towards} established the concept of Epistemic State Specification (ESS) for valid student simulation, arguing that simulators must prioritise the accuracy of the learner's knowledge and reasoning processes over mere dialogue fluency. However, their framework is primarily conceptual and assumes static learner capabilities. Scarlatos et al. \cite{scarlatos2026simulated} benchmarked various simulators, finding that standard prompting strategies suffer heavily from sycophancy bias and fail to maintain consistent misconception-driven belief states. To standardise evaluation across these dimensions, Maurya et al. \cite{maurya2025unifying} propose a unified pedagogical evaluation taxonomy consisting of eight dimensions (including mistake identification, mistake location, and revealing of the answer) and introduce the MRBench math tutoring benchmark. While MRBench provides a rigorous protocol for static dialogue evaluation, it is limited to mathematical domains and does not evaluate adaptive tutoring policies in interactive programming sandboxes.

\subsection{Tutor-Student Agent Interactions}
Frameworks like PETITE \cite{ozdemir2026enhancing} use asymmetric role assignment (a coder student and a helper tutor) to improve code generation success, but they focus purely on task completion rather than pedagogical simulation, lacking explicit models of human cognitive fatigue. SocraticLM \cite{liu2024socraticlm} fine-tunes a model on a Socratic dialogue dataset (SocraTeach) using a Dean-Teacher-Student pipeline, but the tutor agent remains vulnerable to answer leakage when students show persistent confusion. Similarly, the KELE framework \cite{peng2025kele} introduces a structured Socratic teaching methodology via a collaborative `Consultant-Teacher' multi-agent mechanism. While KELE structures dialogue stages using a rule system (SocRule) and demonstrates Socratic performance improvements by fine-tuning on a science-themed dialogue dataset (SocratDataset), its planning is guided by qualitative LLM heuristics. It lacks explicit mathematical formulations of the student's cognitive state dynamics (such as fatigue and friction) and does not mitigate the sycophancy bias of its student simulator. To control dialogue alignment more precisely, Petukhova \& Kochmar \cite{petukhova2025intent} expand Socratic taxonomies (such as MathDial) to 11 fine-grained pedagogical intents, demonstrating that training tutors on specific pedagogical goals improves response alignment. However, their intent-driven mapping is static and does not dynamically adjust to estimated student knowledge states or active misconceptions.

\subsection{Behavioural Simulation \& Persona Modelling}
Cunling Bian \cite{bian2026generative} utilises structured student interviews and multi-expert reflection to simulate learning behaviours (measured by the MSLQ questionnaire) with high accuracy, but does not deploy agents in a live, interactive dialogue environment. Classroom Simulacra \cite{xu2025classroom} incorporates online slide interactions and webcam-driven eye-tracking (GazeRecorder) to capture student engagement, but focuses on classroom-level dynamics rather than dyadic tutor policy optimisation. GPTeach \cite{markel2023gpteach} builds interactive TA training sandboxes using static student personas, but relies on a human teacher in the loop, preventing automated optimisation. 
